% =====================================================================
% O CÁLCULO AUSENTE — versão para submissão ao
% Bolema: Boletim de Educação Matemática (UNESP Rio Claro)
%
% IMPORTANTE: arquivo ANONIMIZADO para avaliação por pares duplo-cega
% (sem nomes/afiliação/ORCID no corpo nem nas propriedades do arquivo).
% Identificação dos autores vai apenas na plataforma ScholarOne e na
% carta de apresentação (calculo-bolema-CARTA-APRESENTACAO.md).
%
% Normas Bolema atendidas: Word/A4; Times New Roman 12; espaço 1,5;
% recuo 1,25 cm; margens sup./esq. 3,0 cm e inf./dir. 2,0 cm; título
% PT 16pt, título secundário 14pt; resumo 100-250 palavras; até 5
% palavras-chave; quadros/tabelas título acima e fonte abaixo, sem
% fios verticais; figuras "Figura N -" com fonte; referências TNR 11;
% ABNT NBR 10520:2023 (autor-data). Limite: 20 páginas.
%
% Compilação: tectonic calculo-bolema-submissao.tex
% =====================================================================
\documentclass[12pt,a4paper,oneside]{article}

\usepackage[T1]{fontenc}
\usepackage[utf8]{inputenc}
\usepackage[brazilian]{babel}
\usepackage{newtxtext}
\usepackage{newtxmath}
\usepackage[a4paper,top=3cm,bottom=2cm,left=3cm,right=2cm]{geometry}
\usepackage{setspace}
\usepackage{indentfirst}
\usepackage{titlesec}
\usepackage{caption}
\usepackage{booktabs}
\usepackage{array}
\usepackage{graphicx}
\usepackage{float}
\usepackage[hidelinks,breaklinks=true]{hyperref}
\usepackage{url}
\usepackage{xurl}

\graphicspath{{figures-calculo/}}

\onehalfspacing
\setlength{\parindent}{1.25cm}
\setlength{\parskip}{0pt}
\setlength{\emergencystretch}{2em}

% --- Seções primárias: TNR 12, negrito, MAIÚSCULAS, justificado ------
\titleformat{\section}{\normalfont\normalsize\bfseries}{\thesection}{0.6em}{\MakeUppercase}
\titleformat{\subsection}{\normalfont\normalsize\bfseries}{\thesubsection}{0.6em}{}
\titleformat{\subsubsection}{\normalfont\normalsize\itshape}{\thesubsubsection}{0.6em}{}
\titlespacing*{\section}{0pt}{14pt}{6pt}
\titlespacing*{\subsection}{0pt}{10pt}{4pt}
\titlespacing*{\subsubsection}{0pt}{8pt}{3pt}

% --- Quadros/Tabelas: título acima, fonte size 10, sem fios verticais
\captionsetup[table]{labelfont=bf,name=Quadro,labelsep=endash,
                     position=above,justification=centering,font=normalsize,skip=4pt}
\captionsetup[figure]{labelfont=bf,name=Figura,labelsep=endash,
                      position=below,justification=centering,font=small,skip=4pt}

% comando para a linha "Fonte:" sob quadros/tabelas
\newcommand{\fonte}[1]{\par\vspace{1pt}{\footnotesize\noindent\hspace*{0pt}Fonte: #1}}

\begin{document}

% =====================================================================
% TÍTULO + RESUMOS  (sem identificação de autoria — avaliação cega)
% =====================================================================
% Título em Português (idioma principal, tamanho 14, máximo 15 palavras)
\begin{center}
{\fontsize{16}{19}\selectfont\bfseries O cálculo ausente: duzentos anos de
currículo, uma comparação internacional e um roteiro pedagógico para o
ensino médio}
\end{center}

\vspace{4pt}

% ---------- RESUMO (PT) ----------
{\centering\textbf{Resumo}\par}\vspace{2pt}
{\setstretch{1.0}\small\noindent
Este artigo conduz uma revisão documental comparada dos currículos
oficiais de matemática do ensino médio, com foco específico na presença
ou ausência de cálculo diferencial e integral antes do ingresso no
ensino superior. Foram analisados documentos curriculares e de avaliação
de sistemas educacionais nacionais em cinco continentes, abrangendo
países de alta, média e baixa renda, além do programa International
Baccalaureate. O achado central é que o Brasil é o único país da amostra
cujo currículo nacional do ensino médio não contempla limites, derivadas
ou integrais. Discutem-se o marco pedagógico que fundamenta a
viabilidade da reintrodução --- ilustrado por uma sequência didática à luz
de Bruner ---, a trajetória
de duzentos anos do currículo brasileiro, do Colégio Pedro II à BNCC, e
indicadores contextuais associados, como o desempenho em avaliações
internacionais e as elevadas taxas de reprovação em Cálculo I nas
engenharias. O achado central é verificável de forma automatizada sobre os
documentos curriculares oficiais. Conclui-se que o vácuo é estrutural, historicamente
documentado e pedagogicamente solucionável, ainda que a inferência
causal estrita entre a ausência curricular e os indicadores observados
permaneça limitada.\par}

\smallskip
\noindent{\small\textbf{Palavras-chave:} Cálculo Diferencial e Integral.
Currículo de Matemática. Ensino Médio. Educação Matemática Comparada.
BNCC.}

\medskip

% Título em Inglês (idioma secundário, tamanho 12)
\begin{center}
{\fontsize{14}{17}\selectfont\bfseries The absent calculus: two hundred years of curriculum, an
international comparison and a pedagogical route for secondary education}
\end{center}
\vspace{2pt}

% ---------- ABSTRACT (EN) ----------
{\centering\textbf{Abstract}\par}\vspace{2pt}
{\setstretch{1.0}\small\noindent
This article presents a comparative documentary review of official
upper-secondary mathematics curricula, focusing specifically on the
presence or absence of differential and integral calculus before
entry into higher education. Curriculum and assessment documents from
national education systems across five continents were analysed,
including high-, middle- and low-income countries, as well as the
International Baccalaureate programme. The central finding is that
Brazil is the only country in the sample whose national upper-secondary
curriculum does not include limits, derivatives or integrals. We discuss
the pedagogical framework supporting the feasibility of reintroduction ---
illustrated by a Bruner-inspired teaching sequence ---,
the two-hundred-year trajectory of the
Brazilian curriculum, from Colégio Pedro II to the BNCC, and associated
contextual indicators, such as performance in international assessments
and high failure rates in university Calculus~I in engineering programmes.
The central finding is verifiable automatically against the official
curricular documents.
We conclude that the gap is structural, historically documented and
pedagogically solvable, although strict causal inference between the
curricular absence and the observed indicators remains limited.\par}

\smallskip
\noindent{\small\textbf{Keywords:} Differential and Integral Calculus.
Mathematics Curriculum. Secondary Education. Comparative Mathematics
Education. BNCC.}

\medskip

% Título em Espanhol (idioma secundário, tamanho 12)
\begin{center}
{\fontsize{14}{17}\selectfont\bfseries El cálculo ausente: doscientos años de currículo, una
comparación internacional y una ruta pedagógica para la enseñanza secundaria}
\end{center}
\vspace{2pt}

% ---------- RESUMEN (ES) ----------
{\centering\textbf{Resumen}\par}\vspace{2pt}
{\setstretch{1.0}\small\noindent
Este artículo realiza una revisión documental comparada de los
currículos oficiales de matemática de la enseñanza secundaria, con foco
específico en la presencia o ausencia de cálculo diferencial e integral
antes del ingreso a la educación superior. Se analizaron documentos
curriculares y de evaluación de sistemas educativos nacionales en cinco
continentes, abarcando países de renta alta, media y baja, además del
programa International Baccalaureate. El hallazgo central es que Brasil
es el único país de la muestra cuyo currículo nacional de secundaria no
contempla límites, derivadas ni integrales. Se discuten el marco
pedagógico que fundamenta la viabilidad de la reintroducción --- ilustrado
por una secuencia didáctica inspirada en Bruner ---, la
trayectoria de doscientos años del currículo
brasileño, del Colégio Pedro II a la BNCC, e indicadores contextuales
asociados, como el desempeño en evaluaciones internacionales y las
elevadas tasas de reprobación en Cálculo~I en las ingenierías. El hallazgo
central es verificable de forma automatizada sobre los documentos
curriculares oficiales. Se
concluye que el vacío es estructural, históricamente documentado y
pedagógicamente solucionable, aunque la inferencia causal estricta
entre la ausencia curricular y los indicadores observados permanezca
limitada.\par}

\smallskip
\noindent{\small\textbf{Palabras clave:} Cálculo Diferencial e Integral.
Currículo de Matemática. Enseñanza Secundaria. Educación Matemática
Comparada. BNCC.}

\bigskip

% =====================================================================
% 1. INTRODUÇÃO
% =====================================================================
\section{Introdução}

A pergunta que orienta este artigo é simples: \emph{em quais países do
mundo um estudante aprende cálculo diferencial e integral antes de
ingressar na universidade?} A resposta surpreende quem está habituado ao
sistema brasileiro. Em praticamente todos os países com tradição
consolidada em educação matemática --- e, como se mostrará, também em
diversos países de renda média e baixa --- o cálculo integra o currículo
do ensino médio dos estudantes que pretendem cursar engenharia, ciências
exatas ou economia. No Brasil, não.

A ausência do cálculo no ensino médio brasileiro é uma anomalia
historicamente conhecida. Geraldo Ávila, no artigo ``O ensino de Cálculo
no 2.\textsuperscript{o} grau'', publicado na \emph{Revista do Professor
de Matemática} em 1991, perguntava diretamente por que não se ensinavam
as ideias básicas do cálculo na escola secundária, argumentando que sua
omissão suprimia ``a componente significativa e certamente a mais
relevante da Matemática'' (Ávila, 1991, p. 1). A questão permanece em
aberto mais de três décadas depois. Wanderley Moura Rezende, em tese de
doutorado defendida na Faculdade de Educação da Universidade de São Paulo
em 2003, formalizou o diagnóstico: a dificuldade dos estudantes
brasileiros em Cálculo~I universitário é \emph{essencialmente
epistemológica}, e tem como raiz a omissão das ideias básicas do cálculo
no ensino que antecede a universidade (Rezende, 2003).

A relevância da pergunta foi recentemente amplificada pelo debate público
sobre o ``apagão de engenheiros''. Segundo o Conselho Federal de
Engenharia e Agronomia, o Brasil pode enfrentar déficit de até um milhão
de profissionais das áreas tecnológicas até 2030, com o número de
matrículas em cursos de engenharia recuando cerca de 30\% entre 2015 e
2024 (CONFEA, 2025). Em paralelo, a disciplina de Cálculo Diferencial e
Integral~I nas engenharias brasileiras apresenta taxas de reprovação
historicamente altas: a Universidade Federal do Rio de Janeiro (UFRJ)
registrou cerca de 70\% de reprovação em período recente (GAZETA DO POVO,
2023), e pesquisa na Universidade Estadual de Campinas (Unicamp)
constatou taxa combinada de reprovação e evasão de até 77,5\% entre 1997
e 2009 (Garzella, 2013).

Há, ainda, uma dimensão de equidade que atravessa todo o argumento. No
Brasil, o cálculo no ensino médio não está totalmente ausente: ele é
estudado, hoje, por estudantes de escolas particulares que oferecem o
programa \emph{International Baccalaureate}, ao passo que a rede pública
--- regida pela BNCC --- não o oferece. A estratificação, portanto,
\emph{já existe}; a questão é se ela deve permanecer invisível na política
curricular. Incorporar o cálculo à base comum tenderia a \emph{reduzir},
não a ampliar, essa desigualdade de acesso --- ponto ao qual se retorna na
discussão (Seção~\ref{sec:discussao}).

O argumento deste artigo não é causal estrito. Não se afirma que a
inclusão de cálculo no ensino médio brasileiro resolveria, por si só, o
conjunto desses problemas. Argumenta-se que (i) o vácuo curricular
existe; (ii) é documentável em fontes primárias; (iii) está registrado na
literatura acadêmica brasileira há mais de três décadas; e (iv) contrasta
com a prática internacional dominante de uma forma que merece ser mapeada
com precisão. Esse mapeamento dialoga com a tradição de estudos
curriculares em Educação Matemática --- da edição temática que a área
dedicou ao currículo de matemática (Pires \emph{et al.}, 2014) aos estudos
comparados entre sistemas nacionais (Gonçalves; Dias; Peralta, 2018) --- e
com a discussão sobre a introdução das ideias do cálculo na educação
básica (Araújo; Avelar, 2022). Diferentemente desses trabalhos, o foco
aqui é estritamente a presença ou ausência do cálculo diferencial e
integral no currículo prescrito, mapeada país a país e verificável de
forma automatizada (Seção~\ref{sec:metodo}).

O restante do artigo organiza-se como segue. A Seção 2 estabelece o
referencial teórico (Bruner, Vygotsky, Rezende e o caso de Singapura). A
Seção 3 descreve a metodologia. A Seção 4 apresenta os resultados e a
discussão: o panorama internacional em países desenvolvidos e em
desenvolvimento, um comparativo por idade, a trajetória histórica
brasileira, as consequências mensuráveis, exemplos de exames e a
discussão dos limites de robustez. A Seção 5 traz as considerações
finais.

% =====================================================================
% 2. REFERENCIAL TEÓRICO
% =====================================================================
\section{Referencial Teórico}

Antes de comparar currículos, é necessário estabelecer o arcabouço
pedagógico que torna a discussão sobre cálculo no ensino médio uma
questão técnica, e não meramente retórica. Três autores são centrais:
Jerome Bruner, Lev Vygotsky e Wanderley Rezende. O caso de Singapura
fornece evidência empírica de que a aplicação sistemática desses
princípios produz resultados internacionais mensuráveis.

\subsection{Bruner e o currículo em espiral}

Jerome Bruner (1915--2016) formulou em 1960 o princípio do \emph{currículo
em espiral}: tópicos matemáticos não devem ser apresentados uma única vez
em sua forma completa, mas revisitados em níveis crescentes de
complexidade ao longo da escolarização, de modo que cada revisita
reconstrói e integra a compreensão prévia (Bruner, 1960). Bruner postulou
ainda três modos de representação do conhecimento que se sucedem na
aprendizagem: \emph{enativo} (ação concreta), \emph{icônico} (imagem
mental) e \emph{simbólico} (notação formal abstrata) (Bruner, 1966).

A consequência prática para o ensino do cálculo é direta: a ideia de
``taxa de variação'' pode ser introduzida no nível enativo (velocidade de
um corpo em queda), depois icônico (gráfico posição-tempo) e, finalmente,
simbólico (definição de derivada via limite). A versão simbólica não
precisa esperar a universidade --- ela é o ponto de chegada de uma
trajetória que pode começar antes. Bruner sintetizou seu princípio em
uma proposição que se tornou clássica: ``qualquer matéria pode ser
ensinada de forma intelectualmente honesta a qualquer criança em qualquer
estágio de desenvolvimento''\footnote{No original: \emph{``any subject
can be taught effectively in some intellectually honest form to any child
at any stage of development''}.} (Bruner, 1960, p. 33).

\subsection{Vygotsky e a zona de desenvolvimento proximal}

Lev Vygotsky (1896--1934) introduziu o conceito de \emph{zona de
desenvolvimento proximal} (ZDP), definida como a distância entre o nível
de desenvolvimento real, determinado pela resolução autônoma de
problemas, e o nível de desenvolvimento potencial, determinado pela
resolução de problemas sob orientação de um adulto ou em colaboração com
pares mais capazes (Vygotsky, 1978, p. 86). A aprendizagem eficaz ocorre
nessa zona. O conceito de \emph{scaffolding} (andaime), formalizado por
Wood, Bruner e Ross, operacionaliza a ZDP: o suporte é fornecido enquanto
necessário e gradualmente retirado conforme o estudante se aproxima de
seu nível potencial (Wood; Bruner; Ross, 1976, p. 90).

O argumento para o cálculo no ensino médio decorre dessa formulação: se
um estudante de 17 anos no Japão, na Alemanha ou em Singapura aprende
derivadas e integrais com \emph{scaffolding} adequado, isso evidencia que
o conteúdo está dentro da ZDP dessa faixa etária. A ausência do conteúdo
no Brasil é decisão curricular, não limitação cognitiva.

\subsection{Rezende e os obstáculos epistemológicos}

A análise mais rigorosa do problema brasileiro foi conduzida por Rezende
(2003). Após análise empírica e teórica das dificuldades de aprendizagem
em Cálculo~I, o autor identificou cinco macro-eixos epistemológicos de
dificuldade: discreto/contínuo; variabilidade/permanência;
finito/infinito; local/global; e sistematização/construção. O diagnóstico
central é que existe \emph{uma fonte comum} dessas dificuldades --- a
omissão das ideias básicas e dos problemas construtores do cálculo no
ensino de matemática que antecede a universidade (Rezende, 2003). Não se
trata, segundo o autor, de falta de pré-requisitos algébricos nem de
inadequação dos métodos universitários, mas do fato de o estudante
encontrar a problemática do cálculo de forma pontual, abstrata e
descontextualizada, sem ter sido preparado para pensar em variação, em
limite e em processos infinitos.

Esse diagnóstico dialoga com a literatura internacional sobre a didática
do cálculo. A dificuldade com a noção de limite é interpretada, na teoria
da \emph{imagem conceitual} e da \emph{definição conceitual}, como um
descompasso entre a intuição do estudante e a formalização (Tall; Vinner,
1981). A derivada, por sua vez, exige a passagem de uma concepção
\emph{operacional} (o processo de calcular taxas) para uma
\emph{estrutural} (a função-derivada como objeto) --- transição que Sfard
(1991) descreve como reificação e que a teoria APOS detalha em termos de
ação--processo--objeto--esquema (Dubinsky; McDonald, 2001). Artigue (1991)
mostra que tais dificuldades são estruturais do ensino da análise, não
idiossincrasias nacionais. A contribuição de Rezende é localizar a
\emph{raiz} brasileira dessas dificuldades na ausência curricular que
antecede a universidade --- o que remete diretamente à noção de
\emph{transposição didática}: o processo pelo qual um saber de referência
(o cálculo, \emph{savoir savant}) é --- ou deixa de ser --- convertido em
saber a ensinar (\emph{savoir enseigné}) no currículo prescrito
(Chevallard, 1991). A ausência documentada neste artigo é, nesses termos,
uma \emph{não-transposição}: uma decisão, e não uma inevitabilidade.

\subsection{O caso de Singapura: Bruner aplicado em escala nacional}
\label{sec:singapura}

Singapura oferece evidência empírica de que a aplicação sistemática dos
princípios de Bruner produz resultados mensuráveis. A partir do início
da década de 1980, o currículo singapurense adotou explicitamente a
sequência \emph{Concrete-Pictorial-Abstract} (CPA) --- operacionalização
direta dos modos enativo, icônico e simbólico de Bruner ---, tal como
documentado pela própria literatura de educação matemática do país
(Leong; Ho; Cheng, 2015). O princípio operacional ficou conhecido como
``menos tópicos, maior profundidade'', com revisitas em espiral: o
conceito de derivada não é apresentado de uma vez como definição-limite,
mas construído ao longo de anos, partindo de razões de variação em
problemas concretos até a diferenciação implícita e as séries de
Maclaurin no \emph{H2 Mathematics} pré-universitário (SEAB, 2025).

O resultado aparece nas avaliações internacionais. Singapura figurou
entre as primeiras posições mundiais em matemática em diversas edições do
TIMSS, incluindo 1995, 1999, 2003 e 2015 (Mullis \emph{et al.}, 2016), e
ocupou a primeira posição global em matemática no PISA 2022, com 41\% de
\emph{top performers} (OECD, 2023). A lição relevante para o Brasil não é
``copiar Singapura'' --- o contexto é distinto ---, mas reconhecer que a
teoria pedagógica que permite ensinar cálculo no ensino médio está
disponível, foi testada em escala nacional por quatro décadas e produz
resultados. A questão deixa de ser ``é possível?'' para tornar-se ``por
que não?''.

\subsection{Do referencial à sala de aula: uma sequência em engenharia didática}
\label{sec:sequencia}

Para que o argumento não seja apenas diagnóstico, convém mostrar como o
referencial se traduz em ação didática. Apresenta-se o \emph{desenho} de
uma sequência --- na tradição da \emph{engenharia didática} (Artigue,
1988), que articula uma análise \emph{a priori} do saber e das
concepções dos estudantes a uma proposta de intervenção posteriormente
validável em sala. O Quadro~\ref{quad:cpa} organiza a introdução tanto da
\emph{derivada} quanto da \emph{integral} no ensino médio brasileiro ---
cobrindo os dois tópicos que o artigo identifica como ausentes --- segundo
os três modos de representação de Bruner (enativo, icônico, simbólico), a
mesma lógica \emph{Concrete-Pictorial-Abstract} operacionalizada por
Singapura. Cada etapa cabe em conteúdos já presentes no currículo
(cinemática, funções, geometria analítica), o que sustenta a viabilidade.

\begin{table}[H]
\centering
\caption{Sequência ilustrativa para introduzir derivada e integral no
ensino médio, segundo os modos de representação de Bruner}
\label{quad:cpa}
\small
\begin{tabular}{@{}p{1.9cm}p{5.6cm}p{4.5cm}@{}}
\toprule
\textbf{Modo} & \textbf{Situação (derivada / integral)} & \textbf{Objeto matemático} \\
\midrule
Enativo \newline (concreto) &
Móvel: medir posição e tempo (velocidade); e recuperar a distância
percorrida a partir da velocidade &
Taxa média $\Delta s/\Delta t$; distância como área acumulada \\
Icônico \newline (gráfico) &
Gráfico posição--tempo (secantes $\to$ tangente); e área sob o gráfico
velocidade--tempo por retângulos &
Tangente como limite de secantes; área como limite de somas de retângulos \\
Simbólico \newline (formal) &
Razão incremental $\to$ limite; e somas de Riemann $\to$ integral definida &
$f'(x)=\lim_{h\to 0}\frac{f(x+h)-f(x)}{h}$ e $\int_a^b f$; o TFC liga as duas \\
\bottomrule
\end{tabular}
\fonte{Elaboração própria, a partir de Bruner (1960, 1966) e da sequência
CPA (Leong; Ho; Cheng, 2015).}
\end{table}

A sequência dialoga diretamente com o diagnóstico de Rezende (2003): ao
partir do concreto e do gráfico antes do símbolo, ela ataca justamente os
eixos epistemológicos que o autor identifica como fontes de dificuldade
--- o par discreto/contínuo (da razão média à taxa instantânea), o
finito/infinito (o processo de limite) e o local/global (a tangente como
comportamento local de uma função). A análise \emph{a priori} é, portanto,
explícita: espera-se que a ancoragem enativa e icônica reduza o
descompasso entre imagem e definição conceitual do limite (Tall; Vinner,
1981) e favoreça a reificação da derivada como objeto (Sfard, 1991) antes
da formalização simbólica. Enquadrados assim, os obstáculos deixam de ser
tratados na universidade, de forma pontual, e passam a ser construídos ao
longo do ensino médio.

Este artigo apresenta o \emph{desenho} e a análise \emph{a priori} da
sequência, não sua validação empírica --- que constitui a etapa seguinte
da engenharia didática. Um protocolo de validação está delineado para
trabalho futuro: aplicação da sequência em turmas do ensino médio, com
análise \emph{a posteriori} contrastando o comportamento observado às
hipóteses aqui declaradas (por exemplo, a evolução da imagem conceitual de
taxa de variação e a passagem da concepção operacional à estrutural da
derivada), em desenho quase-experimental com grupo de comparação. Delimitar
claramente o que é proposta fundamentada e o que ainda demanda evidência de
sala é parte do compromisso metodológico do trabalho.

% =====================================================================
% 3. METODOLOGIA
% =====================================================================
\section{Metodologia}
\label{sec:metodo}

Trata-se de uma pesquisa de abordagem qualitativa, do tipo revisão
documental comparada, de caráter exploratório-descritivo. O corpus é
constituído por documentos curriculares e de avaliação oficiais de cada
sistema educacional --- ministérios da educação, conselhos de currículo e
órgãos de exame nacionais ---, sempre que disponíveis em formato
eletrônico, complementados por relatórios de comparações internacionais
(PISA e TIMSS) e por literatura acadêmica revisada por pares quando havia
divergência ou ambiguidade entre fontes nacionais.

Para cada sistema, catalogaram-se: (i) a estrutura geral do ensino médio
(duração e divisão em trilhas/séries); (ii) o currículo de matemática nos
anos finais (idades de 15 a 18 anos); (iii) a presença ou ausência de
cálculo diferencial (limites, derivadas); (iv) a presença ou ausência de
cálculo integral; (v) o sistema de exame nacional ou de admissão e o peso
atribuído à matemática; e (vi) estatísticas de desempenho disponíveis.

A definição operacional de ``cálculo no currículo'' adotada foi: o tópico
aparece de forma explícita, obrigatória ou eletiva formal, no documento
curricular nacional, com expectativa de avaliação em exame de
saída/admissão. Tópicos eventualmente mencionados em livros didáticos,
mas não previstos no documento curricular nem cobrados em exame nacional,
foram classificados como ``ausentes'' do currículo --- caso reconhecido,
no Brasil, pela própria Sociedade Brasileira de Educação Matemática
(SBEM, 2007).

Para que o achado central não dependa apenas da leitura dos autores, a
catalogação qualitativa é triangulada por uma verificação automatizada e
reprodutível. Um manifesto público (\emph{sources\_calculo\_curricula.json})
registra, para cada sistema, o documento curricular oficial (com URL, órgão
emissor, \emph{hash} SHA-256 e data de captura), e um roteiro
(\emph{audit\_curricula\_keywords.py}) extrai o texto de cada documento
--- inclusive do formato HWP coreano --- e conta ocorrências de termos
inequívocos de cálculo em cada idioma (por exemplo, \emph{cálculo
diferencial}, \emph{derivada de}, \emph{integral definida}; 微分/積分;
미분/적분). Deliberadamente evitam-se termos polissêmicos (``integral'',
``limite'' isolados), que gerariam falso-positivo em português. Sobre os
documentos oficiais dos dez países e do programa IB, o procedimento
confirma a presença de cálculo em todos os sistemas exceto o brasileiro,
cuja BNCC retorna \emph{zero} ocorrências --- tornando o achado verificável
por terceiros em minutos, e não em horas de leitura (dados e código
disponíveis; ver Declarações).

Declaram-se as seguintes limitações metodológicas: (i) documentos
curriculares são \emph{desideratos prescritivos}, não medidas do que de
fato é ensinado em sala, podendo haver descolamento entre currículo
oficial e implementação real em qualquer país; (ii) a análise comparativa
internacional carrega risco de viés do analista, mitigado pela ênfase em
fontes primárias oficiais; (iii) inferências causais entre
presença/ausência de cálculo no ensino médio e indicadores de desempenho
universitário são tratadas como plausíveis, não como demonstradas
(Seção~\ref{sec:discussao}); e (iv) para alguns sistemas
descentralizados (Estados Unidos, Canadá, Alemanha) adotou-se um
documento de referência nacional ou de uma jurisdição representativa, o
que é explicitado caso a caso.

Declara-se, ainda, o uso de assistente de inteligência artificial
(Claude, modelo Opus 4.8, da Anthropic) como ferramenta de apoio em
todas as etapas do trabalho --- organização e revisão das fontes
curriculares, redação e edição do texto e geração do código das figuras.
Todo o conteúdo foi revisado e verificado pelos autores, que assumem
integral responsabilidade pública pelo texto (ver Declarações).

% =====================================================================
% 4. RESULTADOS E DISCUSSÃO
% =====================================================================
\section{Análises e Resultados}

\subsection{Panorama internacional: países de alta renda}

O detalhamento por país consta da Figura~\ref{fig:heatmap} e do
repositório de reprodutibilidade; aqui, sintetiza-se o padrão por região.

\paragraph{Ásia.} Nos grandes sistemas asiáticos, o cálculo integra a
trilha de ensino médio voltada a ciências e engenharia: o \emph{Math~III}
japonês cobre derivadas avançadas e integração (JASSO, 2025); o currículo
chinês de 2017 inclui derivadas e, na trilha avançada do Gaokao, a
integral definida e o Teorema Fundamental do Cálculo (China, 2020); a
Coreia do Sul mantém limites e continuidade em disciplina obrigatória e o
cálculo como eletiva típica dos candidatos STEM (Coreia do Sul, 2015); e o
\emph{H2 Mathematics} de Singapura cobre diferenciação, integração e
equações diferenciais (SEAB, 2025).

\paragraph{Europa.} O padrão se repete. O cálculo é compulsório na trilha
científica no \emph{Abitur} alemão (Alemanha, 2012), no \emph{Liceo
Scientifico} italiano (Itália, 2010), na \emph{Matemáticas~II} espanhola
(Espanha, 2022) e no nível \emph{Profile} russo (Rússia, 2024); e integra
opções de trilha avançada tipicamente escolhidas por candidatos a STEM na
\emph{Spécialité} francesa (França, 2019), no \emph{A-Level} britânico
(Reino Unido, 2016) e na ``matemática longa'' finlandesa (Finlândia, 2019).

\paragraph{Américas, Oceania e programa internacional.} Nos Estados
Unidos, sem currículo nacional, o \emph{AP Calculus} é a referência de
padrão, em caráter eletivo (College Board, 2020); o Canadá
(\emph{Calculus and Vectors}, Ontário) e a Austrália (\emph{Mathematical
Methods}/\emph{Specialist}) o mantêm na trilha STEM (Canadá, 2007;
Austrália, 2015). Em Israel, o nível de cinco unidades dedica cerca de um
quarto da carga ao cálculo (Kontorovich, 2021). O \emph{International
Baccalaureate} (\emph{Analysis and Approaches} HL), presente em escolas
particulares brasileiras, inclui limites, diferenciação, integração e
equações diferenciais, com a maior carga do programa (IBO, 2019) --- o
ponto de equidade já assinalado.

\subsection{Panorama internacional: países de renda média e baixa}
\label{sec:desenvolvimento}

Uma objeção previsível é que a comparação anterior privilegiaria países
ricos, configurando um artefato de renda. Para testá-la, examinaram-se
sistemas educacionais de renda média e baixa, inclusive de países com
Índice de Desenvolvimento Humano (IDH) \emph{inferior} ao do Brasil
(0,760 em 2022). O resultado é direto: a presença do cálculo no ensino
médio não acompanha a renda nacional.

Entre as grandes nações em desenvolvimento, incluem cálculo diferencial e
integral no ensino médio a Índia --- onde o bloco responde por 35 dos 80
pontos do exame nacional (CBSE, 2023) ---, o Vietnã, como conteúdo nuclear
(Vietnã, 2018), o Paquistão, cujo livro oficial da 12.\textsuperscript{a}
série se intitula \emph{Calculus and Analytic Geometry} (Paquistão, 2023),
Bangladesh, na trilha de ciências (Bangladesh, 2022), e a Nigéria, via
\emph{Further Mathematics} do WASSCE (WAEC, 2023) --- todos com IDH
\emph{inferior} ao brasileiro em vários casos. Entre os pares
latino-americanos, o México torna obrigatório o \emph{Cálculo Diferencial}
(México, 2023), o Chile oferece o eletivo ``Límites, derivadas e
integrales'' (Chile, 2021) e a Argentina prevê limites e derivadas nos
desenhos curriculares provinciais (Argentina, 2011).

Casos intermediários reforçam --- e não enfraquecem --- a distinção
brasileira: África do Sul (África do Sul, 2011), Turquia (Turquia, 2018) e
Colômbia (Colômbia, 2006) mantêm o cálculo \emph{diferencial}, ainda que
restrinjam ou venham removendo a integral. Nenhum sistema da amostra
exclui o cálculo \emph{integralmente} do currículo nacional --- como faz a
BNCC brasileira.

Esse achado obriga a uma formulação precisa do argumento central, adotada
no restante do artigo: \textbf{o Brasil é o único país da amostra cujo
currículo nacional do ensino médio exclui, simultânea e integralmente, o
cálculo diferencial \emph{e} o cálculo integral}. A Figura~\ref{fig:hdi}
torna visual esse resultado: ao distribuir os países de renda média e
baixa examinados por IDH e por grau de presença do cálculo, evidencia-se
que vários países com IDH \emph{inferior} ao brasileiro (Índia 0,644;
Paquistão 0,540; Nigéria 0,548; Vietnã 0,726) oferecem cálculo completo,
enquanto o Brasil (IDH 0,760) é o único ponto na categoria ``ausente''.

\begin{figure}[H]
\centering
\includegraphics[width=0.80\textwidth]{fig03_hdi_calculo.pdf}
\caption{Índice de Desenvolvimento Humano (2022) versus presença de
cálculo no ensino médio, para os países de renda média e baixa examinados
e o Brasil. Vários países com IDH \emph{inferior} ao
brasileiro (Índia, Vietnã, Paquistão, Nigéria) oferecem cálculo completo;
o Brasil (vermelho) é o único ponto na categoria ``ausente''.}
\label{fig:hdi}
\end{figure}

\subsection{Síntese quantitativa}
\label{sec:sintese}

A Figura~\ref{fig:heatmap} sintetiza a presença de cálculo no ensino médio
dos sistemas analisados, célula a célula (país $\times$ tópico), com o
Brasil como única linha integralmente ausente. Os documentos oficiais de
cada sistema, o status curricular codificado e o desempenho no PISA 2022
por país constam do repositório de reprodutibilidade (ver Declarações).

\begin{figure}[H]
\centering
\includegraphics[width=0.95\textwidth]{fig01_heatmap_calculo.pdf}
\caption{Cálculo no ensino médio em dezessete sistemas nacionais e no
programa IB (subconjunto que preserva a variação relevante; a matriz
completa consta do repositório de auditoria). Cada célula indica se o
tópico (limites, derivadas, integrais, equações diferenciais, séries,
geometria diferencial) é obrigatório, eletivo ou ausente no currículo
nacional para ingresso em cursos STEM. O Brasil (BNCC) é a única linha
integralmente ausente.}
\label{fig:heatmap}
\end{figure}

Uma verificação estatística acompanha esse panorama. Codificando o status
curricular de forma ordinal (obrigatório~$=2$; eletivo/opção de trilha
avançada~$=1$; ausente~$=0$ --- codificação documentada no script
\emph{compute\_pisa\_correlations.py}) e correlacionando-o com o desempenho
no PISA 2022 dos quinze sistemas com dado disponível, o coeficiente de correlação
de postos de Spearman é nulo ($\rho = -0{,}08$; IC~95\% por \emph{bootstrap}
$[-0{,}67,\, 0{,}51]$; teste de permutação unilateral $p \approx 0{,}61$,
hipótese de gradiente positivo). A não-significância é estável sob
codificações ordinais alternativas do status curricular. Isto é: entre os
países que \emph{têm} cálculo, o grau de obrigatoriedade não prediz o desempenho ---
resultado coerente com a ressalva de que o PISA é um \emph{proxy} invertido
(Seção~\ref{sec:discussao}). O que é robusto não é um gradiente
dose-resposta, e sim a posição do Brasil como \emph{outlier}: seu escore
(379) está a cerca de $3{,}5$ desvios-padrão abaixo da média dos sistemas
com cálculo (492). O contraste ausente-\emph{versus}-presente também é
marcante (correlação ponto-bisserial $r = 0{,}68$; $p < 0{,}01$), embora,
com o Brasil como único caso ``ausente'', esse coeficiente seja uma
reparametrização do mesmo escore-$z$, não uma evidência independente. O
código que reproduz estas estatísticas é público (ver Declarações). Os
dados sustentam, portanto, a \emph{singularidade curricular} do Brasil ---
não uma relação causal entre cálculo no ensino médio e desempenho aos 15 anos.

\subsection{Comparativo por idade: o que se aprende em cada faixa etária}

Tomado um estudante qualquer dos países comparados, em que idade ele
encontra qual conteúdo? Já aos 16 anos o distanciamento se abre: na maioria
dos sistemas, o estudante realiza cálculo diferencial em algum nível
(diferenciação em Singapura e na Alemanha; derivadas e aplicações na China;
limites na Coreia do Sul), enquanto no Brasil o conteúdo é análise
combinatória e geometria. Aos 17 anos, a divergência se completa: Japão
(\emph{Math~III}), China, Singapura, Alemanha, França, Itália, Reino Unido
e Rússia encerram o ensino médio com integração e equações diferenciais
--- o equivalente ao Cálculo~I universitário brasileiro ---, ao passo que a
matemática do 3.\textsuperscript{o} ano brasileiro cobre tópicos de
pré-cálculo (geometria analítica, polinômios) que \emph{seriam}
pré-requisitos do cálculo, sem nunca chegar a ele. Bruner (1960) descreveu
esse padrão como antipadrão pedagógico: apresentar um conceito complexo em
sua forma final, sem revisitas progressivas, tende a produzir compreensão
superficial e baixa transferência.

\subsection{O caso brasileiro: duzentos anos de currículo}

O caso brasileiro só se compreende em sua trajetória histórica. Ao
contrário do que sugere o senso comum, a ausência do cálculo no ensino
médio não é ``natural'' nem ``definitiva'': resulta de decisões
curriculares sucessivas, algumas das quais incluíram o cálculo no
programa antes de removê-lo.

O Colégio Pedro II (decreto de 1837), primeira instituição secundária
organizada do país, tinha currículo humanístico com matemática elementar,
sem cálculo (Brasil, 1837). A primeira inclusão formal veio com a Reforma
Benjamin Constant: o Decreto n.\textsuperscript{o} 981, de 1890, inspirado
na hierarquia das ciências de Comte, previu no 3.\textsuperscript{o} ano
``cálculo diferencial e integral, limitado ao conhecimento das teorias
rigorosamente indispensáveis ao estudo da mecânica geral'' (Brasil, 1890).
O cálculo foi excluído já em 1899 e a exclusão consolidada pelo Decreto
n.\textsuperscript{o} 3.890, de 1901 (Silva, 2023); desde então não
retornou ao currículo oficial.

Essa janela 1890--1901 é a peça de evidência mais próxima de um
experimento natural disponível neste desenho: a remoção foi \emph{exógena}
ao desempenho dos estudantes --- seguiu-se à morte prematura de Benjamin
Constant (janeiro de 1891) e ao realinhamento político, não a uma
avaliação técnica de ineficácia (Silva, 2023) ---, ainda que não existam
medidas de desempenho contemporâneas que permitam estimar seu efeito
(Seção~\ref{sec:discussao}).

A Reforma Capanema (Decreto-Lei n.\textsuperscript{o} 4.244/1942),
incorporando propostas de Euclides Roxo, recuperou o \emph{estudo das
variações e das derivadas} no curso Científico, mas não o cálculo integral
(Dassie, 2001). O Movimento da Matemática Moderna (anos 1960), disseminado
por Sangiorgi e pelo GEEM, enfatizou conjuntos e estruturas em detrimento
do cálculo, representando novo recuo (Búrigo, 1989; Soares, 2001). Os PCN
(1997/2000) oficializaram estatística e probabilidade; a BNCC (2018)
organiza a área por competências e \emph{não} inclui limites, derivadas ou
integrais, com estatística apenas descritiva (Brasil, 2018). A própria
SBEM reconhece o problema:

\begin{quote}
\small
``Alguns livros didáticos do Ensino Médio apresentam tópicos de Cálculo
Diferencial e Integral, como limite, derivada e integral, mas na maioria
das vezes não são ensinados sob o pretexto de serem difíceis e impróprios
a esse segmento da educação [\ldots]. Assim sendo, o Cálculo faz parte do
livro didático, mas não do currículo do Ensino Médio.'' (SBEM, 2007).
\end{quote}

O Novo Ensino Médio (Lei n.\textsuperscript{o} 13.415/2017, reformulada
pela Lei n.\textsuperscript{o} 14.945/2024) instituiu itinerários
formativos, entre eles ``Matemática e suas Tecnologias'', que
\emph{permitem}, mas não \emph{exigem}, a oferta de cálculo, sem que isso
o tenha incorporado ao currículo comum obrigatório (Brasil, 2017; BRASIL,
2024).

\subsection{Consequências mensuráveis}

\subsubsection{Desempenho no PISA (2003--2022)}

Uma ressalva de identificação precede os dados. O PISA avalia estudantes
de 15 anos, idade em que o cálculo ainda \emph{não} foi introduzido em
nenhum dos países da amostra; logo, o desempenho brasileiro mede sobretudo
a qualidade do ensino fundamental e as condições socioeconômicas, e não o
efeito específico da ausência de cálculo no ensino médio. Os números a
seguir são, portanto, contexto histórico-regional --- não teste do canal
causal discutido na Seção~\ref{sec:discussao}.

O Brasil melhorou modestamente entre 2003 e 2009 e estagnou em torno de
380 pontos em matemática desde então, mantendo-se cerca de 90 a 145
pontos abaixo da média da OECD ao longo do período --- de $-$143 pontos em
2003 a $-$93 em 2022 (OECD, 2023; Figura~\ref{fig:pisa}). Apenas 1\% dos
jovens brasileiros qualificam-se como \emph{top performers} em matemática,
contra 41\% em Singapura e 9\% na média da OECD (OECD, 2023). O baixo
desempenho persiste mesmo na comparação regional: o Brasil figura na base
do grupo latino-americano, abaixo de Chile, Uruguai, México, Peru e
Colômbia, e em patamar próximo ao da Argentina.

\begin{figure}[H]
\centering
\includegraphics[width=0.82\textwidth]{fig02_pisa_timeline.pdf}
\caption{Desempenho em Matemática no PISA, 2003--2022, em vinte sistemas
educacionais. O Brasil (vermelho) oscila em torno de 380 pontos,
situando-se na base mesmo entre os pares latino-americanos (laranja); a
média da OECD (azul tracejada) e os líderes mundiais e economias
desenvolvidas (cinza) situam-se muito acima. Como o PISA avalia estudantes
de 15 anos, a figura é contexto histórico-regional, não teste do canal
causal. O resultado da Argentina em
2015 foi omitido por incomparabilidade reconhecida pela OECD.}
\label{fig:pisa}
\end{figure}

\subsubsection{Reprovação em Cálculo~I}

A literatura é convergente quanto ao fato de o Cálculo~I ser, isoladamente,
a disciplina que mais reprova nas engenharias brasileiras, embora não
exista um índice nacional único e auditado. Estudos institucionais
reportam taxas elevadas e dispersas: cerca de 70\% na UFRJ (GAZETA DO
POVO, 2023); até 77,5\% de reprovação somada à evasão na Unicamp, entre
1997 e 2009 (Garzella, 2013); e ampla variação entre instituições em
revisões da literatura de educação em engenharia (Macambira; Athayde,
2014). A análise epistemológica de Rezende (2003) interpreta o fenômeno
como estrutural, e não conjuntural.

Em nome do rigor, registre-se que altas taxas de reprovação em cálculo
introdutório não são exclusividade brasileira: nos Estados Unidos, as
taxas de reprovação e abandono (\emph{D/F/W}) em \emph{Calculus~I}
universitário situam-se com frequência entre 25\% e 40\%, mesmo com o
cálculo presente no ensino médio via \emph{AP} (Bressoud, 2015). O
argumento deste artigo, portanto, não é que a ausência de cálculo no
ensino médio \emph{explique} a reprovação --- que tem múltiplas causas
--- mas que o Brasil enfrenta esse filtro \emph{sem} a preparação prévia
que os demais sistemas oferecem.

\subsubsection{O ``apagão de engenheiros''}

O CONFEA projeta déficit de até um milhão de profissionais das áreas
tecnológicas até 2030, associando-o à queda de matrículas e à fragilidade
da formação matemática de base (CONFEA, 2025). A causalidade entre
formação básica e déficit profissional não é direta nem única, mas a
literatura identifica o ciclo básico --- e o Cálculo~I em particular ---
como principal filtro de evasão nas engenharias.

\subsection{Exemplos comparativos: o que se cobra em cada exame?}

Para tornar concreta a discussão, apresentam-se exemplos típicos de
problemas de exames nacionais ao final do ensino médio. Os enunciados
estrangeiros exigem cálculo formal; o brasileiro, aritmética --- na mesma
faixa etária.

\smallskip
\noindent\fbox{\parbox{0.95\textwidth}{\small
\textbf{Gaokao (China):} seja $f(x) = x^3 - 3ax + 2$; discuta a
monotonicidade de $f$, determine $a$ para que o mínimo seja $0$ e o número
de raízes de $f(x)=1$ --- \emph{exige derivar e analisar variação}
(análogos formais constam do Abitur alemão e do AP Calculus BC).\\[3pt]
\textbf{ENEM (Brasil):} o número de bactérias dobra a cada hora; a partir
de $100$, quantas haverá após 5 horas? --- \emph{crescimento exponencial
em forma discreta, sem interface com a derivada nem com a taxa instantânea}.}}

\smallskip
Na mesma faixa etária, um exame cobra cálculo formal; o outro, aritmética
--- contraste que ilustra, de modo tangível, o vácuo curricular documentado.

\subsection{Discussão: níveis de robustez e contra-argumentos}
\label{sec:discussao}

A análise permite distinguir três níveis de afirmação, com graus
distintos de robustez. \textbf{Nível 1 --- fato curricular}: o Brasil é o
único país da amostra cujo currículo nacional do ensino médio não inclui
cálculo (robustez alta). \textbf{Nível 2 --- padrões associados}: o Brasil
é um \emph{outlier} --- desempenho substancialmente inferior em avaliações
internacionais e taxas elevadas de reprovação em Cálculo~I (robustez alta
para a \emph{posição} do Brasil). Note-se, porém, que \emph{não} há
gradiente dose-resposta detectável entre o grau de obrigatoriedade do
cálculo e o desempenho no PISA ($\rho = -0{,}08$; Seção~\ref{sec:sintese}) ---
o que é coerente com o fato de o PISA medir alunos de 15 anos.
\textbf{Nível 3 --- inferência causal}: que a ausência de cálculo no
ensino médio \emph{contribua} para esses indicadores é hipótese plausível,
mas não demonstrada (robustez baixa; há fatores confundidores, como
qualidade geral do ensino básico, condições socioeconômicas e formação
docente).

Em fidelidade ao princípio acadêmico de considerar os argumentos
contrários em sua versão mais forte, registram-se as principais objeções:
(C1)~``a base é fraca demais'' --- objeção parcialmente válida, mas
endereçável por sequenciamento em espiral; (C2)~``não há professores
qualificados'' --- objeção empiricamente forte, que demanda investimento
em formação docente; (C3)~``não cabe na carga horária'' --- a comparação
internacional mostra cargas semanais maiores em vários sistemas,
sugerindo reorganização de prioridades; e (C4)~``aumentaria a
desigualdade'' --- objeção politicamente relevante, mas tecnicamente
reversível, pois a estratificação \emph{já} existe (estudantes de IB em
escolas particulares brasileiras estudam cálculo, ao passo que a rede
pública não o oferece), de modo que incorporá-lo à BNCC tenderia a
\emph{reduzir} a desigualdade.

% =====================================================================
% 5. CONSIDERAÇÕES FINAIS
% =====================================================================
\section{Considerações Finais}

Este artigo documentou, com base em fontes curriculares oficiais de
sistemas educacionais de cinco continentes, que o Brasil é o único da
amostra cujo currículo nacional do ensino médio não inclui cálculo
diferencial e integral. O fato é verificável e historicamente conhecido:
em duzentos anos de trajetória educacional, o cálculo esteve presente no
secundário brasileiro apenas em um período breve --- entre a Reforma
Benjamin Constant (1890) e sua reversão (1899--1901) ---, com a derivada
brevemente recuperada, em caráter restrito, pela Reforma Capanema (1942).

O arcabouço pedagógico que tornaria viável a reintrodução está
consolidado --- o currículo em espiral de Bruner, a zona de
desenvolvimento proximal de Vygotsky e o diagnóstico epistemológico de
Rezende ---, e o caso de Singapura evidencia, em escala nacional, que sua
aplicação sistemática produz resultados. A inferência causal estrita
entre o vácuo curricular e os indicadores mensuráveis permanece difícil
de estabelecer; a posição atípica do Brasil, contudo, é robusta, e a
literatura brasileira converge, desde Ávila (1991), para a hipótese de que
a intervenção curricular seria parte da solução.

Mais do que prescrever uma reforma --- tarefa de longa maturação política
e técnica ---, o artigo pretende oferecer um mapeamento preciso e
replicável de um contraste internacional pouco quantificado na literatura
nacional. Pesquisas futuras poderão (i) ampliar a amostra de países de
renda média e baixa; (ii) realizar análise de séries temporais
interrompidas sobre coortes históricas de ingresso em engenharia; e
(iii) desenhar e avaliar experimentalmente sequências didáticas de
introdução ao cálculo no ensino médio brasileiro. O Brasil é exceção
curricular; não precisa continuar sendo.

% ---------- AGRADECIMENTOS / DECLARAÇÕES (anonimizadas) --------------
\section*{Agradecimentos}
\textit{[Suprimido para avaliação por pares duplo-cega. Será inserido na
versão final, conforme política da revista.]}

\section*{Declarações}
\noindent\textbf{Conflito de interesse.} Os autores declaram não haver
conflito de interesse. \textbf{Financiamento.} A pesquisa não recebeu
financiamento externo. \textbf{Contribuição.} Todos os autores
contribuíram para a concepção do estudo, a curadoria das fontes e a
redação, tendo aprovado a versão submetida. \textbf{Uso de inteligência
artificial.} Os autores declaram o uso do assistente de IA Claude
(modelo Opus 4.8, da Anthropic) como ferramenta de apoio em todas as
etapas de elaboração deste trabalho --- organização e revisão das fontes,
redação e edição do texto e geração do código das figuras. Nenhum
conteúdo foi publicado sem revisão e verificação humanas; os autores
assumem integral responsabilidade pública pelo texto.
\textbf{Disponibilidade de dados e código.} Todos os artefatos que
reproduzem os achados são públicos, em repositório de acesso aberto: o
manifesto das fontes curriculares oficiais (\emph{sources\_calculo\_curricula.json},
com URL, órgão emissor, \emph{hash} SHA-256 e data de captura); o roteiro
de auditoria que reproduz o achado curricular (\emph{audit\_curricula\_keywords.py});
o script que reproduz as correlações com o PISA
(\emph{compute\_pisa\_correlations.py}); e as dependências fixadas
(\emph{requirements-audit.txt}). Para a avaliação por pares, os artefatos
estão disponíveis em espelho anônimo (\emph{double-blind}):
\url{https://anonymous.4open.science/r/calculo-ausente-repro/}. A URL
definitiva do repositório será informada na versão final.

% =====================================================================
% REFERÊNCIAS  (ABNT NBR 6023; chamadas autor-data NBR 10520:2023)
% =====================================================================
\newpage
\section*{Referências}
\addcontentsline{toc}{section}{Referências}

\begingroup
\small
\setstretch{1.0}
\setlength{\parindent}{0pt}
\raggedright
\newcommand{\refitem}[1]{\par #1\par\vspace{1pt}}

\refitem{ÁFRICA DO SUL. Department of Basic Education. \textbf{Curriculum
and Assessment Policy Statement (CAPS), FET Phase, Grades 10--12:
Mathematics}. Pretória: DBE, 2011.}

\refitem{ALEMANHA. Kultusministerkonferenz (KMK). \textbf{Bildungsstandards
im Fach Mathematik für die Allgemeine Hochschulreife}. Berlim: KMK, 2012.
Disponível em:
\url{https://www.kmk.org/fileadmin/veroeffentlichungen_beschluesse/2012/2012_10_18-Bildungsstandards-Mathe-Abi.pdf}.
Acesso em: 6 jun. 2026.}

\refitem{ARAÚJO, Jussara de Loiola; AVELAR, Petrina Rúbria Nogueira.
Modelagem Matemática e o desenvolvimento do pensamento integral.
\textbf{Bolema: Boletim de Educação Matemática}, Rio Claro, v. 36, n. 74,
2022.}

\refitem{ARGENTINA. Dirección General de Cultura y Educación (Provincia
de Buenos Aires). \textbf{Diseño Curricular para la Educación Secundaria
--- Matemática, Ciclo Superior, 6.\textsuperscript{o} año}. La Plata:
DGCyE, 2011.}

\refitem{ARTIGUE, Michèle. Ingénierie didactique. \textbf{Recherches en
Didactique des Mathématiques}, Grenoble, v. 9, n. 3, p. 281--308, 1988.}

\refitem{ARTIGUE, Michèle. Analysis. \emph{In}: TALL, David (org.).
\textbf{Advanced Mathematical Thinking}. Dordrecht: Kluwer, 1991. p. 167--198.}

\refitem{AUSTRÁLIA. Australian Curriculum, Assessment and Reporting
Authority (ACARA). \textbf{The Australian Curriculum: Senior Secondary ---
Mathematics} (v.~8.4). Sydney: ACARA, 2015. Disponível em:
\url{https://www.australiancurriculum.edu.au/senior-secondary-curriculum/mathematics/}.
Acesso em: 6 jun. 2026.}

\refitem{ÁVILA, Geraldo. O ensino de Cálculo no 2.\textsuperscript{o}
grau. \textbf{Revista do Professor de Matemática}, Rio de Janeiro, n. 18,
p. 1--9, 1991.}

\refitem{BANGLADESH. National Curriculum and Textbook Board (NCTB).
\textbf{Higher Secondary (XI--XII) Higher Mathematics Syllabus}. Daca:
NCTB, 2022.}

\refitem{BRASIL. \textbf{Decreto de 2 de dezembro de 1837}. Cria o
Colégio de instrução secundária (Colégio Pedro II). Rio de Janeiro, 1837.}

\refitem{BRASIL. \textbf{Decreto n.\textsuperscript{o} 981, de 8 de
novembro de 1890}. Aprova o Regulamento da Instrução Primária e Secundária
do Distrito Federal (Reforma Benjamin Constant). Rio de Janeiro, 1890.}

\refitem{BRASIL. \textbf{Decreto n.\textsuperscript{o} 3.890, de
1.\textsuperscript{o} de janeiro de 1901}. Código dos Institutos Oficiais
de Ensino Superior e Secundário (Reforma Epitácio Pessoa). Rio de
Janeiro, 1901.}

\refitem{BRASIL. \textbf{Decreto-Lei n.\textsuperscript{o} 4.244, de 9 de
abril de 1942}. Lei Orgânica do Ensino Secundário (Reforma Capanema). Rio
de Janeiro, 1942.}

\refitem{BRASIL. \textbf{Lei n.\textsuperscript{o} 13.415, de 16 de
fevereiro de 2017}. Altera as Leis de Diretrizes e Bases e institui a
Política de Fomento ao Ensino Médio em Tempo Integral. Brasília, 2017.}

\refitem{BRASIL. Ministério da Educação. \textbf{Base Nacional Comum
Curricular}. Brasília: MEC, 2018.}

\refitem{BRASIL. \textbf{Lei n.\textsuperscript{o} 14.945, de 31 de julho
de 2024}. Altera a Lei n.\textsuperscript{o} 9.394/1996 quanto à
organização do ensino médio. Brasília, 2024.}

\refitem{BRESSOUD, David M.; MESA, Vilma; RASMUSSEN, Chris (org.).
\textbf{Insights and recommendations from the MAA National Study of
College Calculus}. Washington, DC: MAA Press, 2015.}

\refitem{BRUNER, Jerome S. \textbf{The process of education}. Cambridge,
MA: Harvard University Press, 1960.}

\refitem{BRUNER, Jerome S. \textbf{Toward a theory of instruction}.
Cambridge, MA: Belknap Press of Harvard University Press, 1966.}

\refitem{BÚRIGO, Elisabete Zardo. \textbf{Movimento da Matemática Moderna
no Brasil}: estudo da ação e do pensamento de educadores matemáticos nos
anos 60. 1989. Dissertação (Mestrado em Educação) --- Universidade Federal
do Rio Grande do Sul, Porto Alegre, 1989.}

\refitem{CANADÁ. Ministry of Education (Ontario). \textbf{The Ontario
Curriculum, Grades 11 and 12: Mathematics} (revised). Toronto: Queen's
Printer for Ontario, 2007. Disponível em:
\url{https://www.edu.gov.on.ca/eng/curriculum/secondary/math1112currb.pdf}.
Acesso em: 6 jun. 2026.}

\refitem{CBSE --- Central Board of Secondary Education.
\textbf{Mathematics (041) Syllabus, Classes XI--XII (2023--24)}. Nova
Délhi: CBSE, 2023.}

\refitem{CHEVALLARD, Yves. \textbf{La transposition didactique}: du savoir
savant au savoir enseigné. 2. ed. Grenoble: La Pensée Sauvage, 1991.}

\refitem{CHILE. Ministerio de Educación (Unidad de Currículum y
Evaluación). \textbf{Programa de Estudio 3.\textsuperscript{o} o
4.\textsuperscript{o} Medio, Formación Diferenciada --- Matemática:
Límites, Derivadas e Integrales}. Santiago: MINEDUC, 2021.}

\refitem{CHINA. Ministério da Educação. \textbf{Currículo nacional de
matemática do ensino médio} [Pu tong gao zhong shu xue ke cheng biao
zhun], ed. 2017, rev. 2020. Pequim: People's Education Press, 2020.}

\refitem{COLLEGE BOARD. \textbf{AP Calculus AB and AP Calculus BC course
and exam description}. Nova York: College Board, 2020. Disponível em:
\url{https://apcentral.collegeboard.org/media/pdf/ap-calculus-ab-and-bc-course-and-exam-description.pdf}.
Acesso em: 6 jun. 2026.}

\refitem{COLÔMBIA. Ministerio de Educación Nacional. \textbf{Estándares
Básicos de Competencias en Matemáticas}. Bogotá: MEN, 2006.}

\refitem{CONFEA. \textbf{Confea alerta para déficit de engenheiros e
defende foco na matemática}. Brasília: Conselho Federal de Engenharia e
Agronomia, 2025. Disponível em:
\url{https://www.confea.org.br/confea-alerta-para-deficit-de-engenheiros-e-defende-foco-na-matematica}.
Acesso em: 6 jun. 2026.}

\refitem{COREIA DO SUL. Ministry of Education. \textbf{Mathematics
Curriculum} (Notification n. 2015-74, Annex 8). Sejong: Ministry of
Education, 2015.}

\refitem{DASSIE, Bruno Alves. \textbf{A matemática do curso secundário na
Reforma Gustavo Capanema}. 2001. Dissertação (Mestrado) --- Pontifícia
Universidade Católica do Rio de Janeiro, Rio de Janeiro, 2001.}

\refitem{DUBINSKY, Ed; McDONALD, Michael A. APOS: a constructivist theory of
learning in undergraduate mathematics education research. \emph{In}: HOLTON,
Derek (org.). \textbf{The Teaching and Learning of Mathematics at University
Level}: an ICMI Study. Dordrecht: Kluwer, 2001. p. 275--282.}

\refitem{ESPANHA. Ministerio de Educación y Formación Profesional.
\textbf{Real Decreto 243/2022}, de 5 de abril, por el que se establecen
la ordenación y las enseñanzas mínimas del Bachillerato. Boletín Oficial
del Estado, Madrid, n. 82, 6 abr. 2022.}

\refitem{FINLÂNDIA. Finnish National Agency for Education (EDUFI).
\textbf{National core curriculum for general upper secondary education
2019}. Helsinki: EDUFI, 2019.}

\refitem{FRANÇA. Ministère de l'Éducation Nationale et de la Jeunesse.
\textbf{Programme de spécialité de mathématiques de la classe terminale
de la voie générale}. Bulletin officiel spécial n. 1, 22 jan. 2019.
Paris: MENJ, 2019.}

\refitem{GARZELLA, Fabiana Colombo. \textbf{A disciplina de Cálculo I}:
análise das relações entre as práticas pedagógicas do professor e seus
impactos nos alunos. 2013. Tese (Doutorado em Educação) --- Universidade
Estadual de Campinas, Campinas, 2013.}

\refitem{GAZETA DO POVO. \textbf{UFRJ vai mudar currículo de cálculo para
tentar evitar evasão}. Curitiba, 2023.}

\refitem{GONÇALVES, Harryson Júnio Lessa; DIAS, Ana Lúcia Braz; PERALTA,
Deise Aparecida. Estudo comparativo sobre o ensino de Matemática em
currículos de educação profissional técnica: Brasil e Estados Unidos.
\textbf{Bolema: Boletim de Educação Matemática}, Rio Claro, v. 32, n. 60,
p. 246--265, 2018.}

\refitem{INTERNATIONAL BACCALAUREATE ORGANIZATION (IBO). \textbf{Diploma
Programme --- Mathematics: Analysis and Approaches guide} (first
assessment 2021). Cardiff: IBO, 2019.}

\refitem{ITÁLIA. Ministero dell'Istruzione, dell'Università e della
Ricerca. \textbf{Indicazioni nazionali riguardanti gli obiettivi
specifici di apprendimento per i percorsi liceali}. Decreto
Interministeriale n. 211, 7 ott. 2010. Roma: MIUR, 2010.}

\refitem{JAPAN STUDENT SERVICES ORGANIZATION (JASSO). \textbf{EJU syllabus
for mathematics}. Tóquio: JASSO, 2025. Disponível em:
\url{https://www.jasso.go.jp/en/ryugaku/eju/examinee/syllabus/mathematics.html}.
Acesso em: 6 jun. 2026.}

\refitem{KONTOROVICH, Igor'. Research as a resource in a high-school
calculus curriculum. \textbf{ZDM --- Mathematics Education}, v. 53, p.
429--440, 2021. DOI: 10.1007/s11858-021-01236-3.}

\refitem{LEONG, Yew Hoong; HO, Weng Kin; CHENG, Lu Pien.
Concrete-Pictorial-Abstract: surveying its origins and charting its
future. \textbf{The Mathematics Educator}, Singapura, v. 16, n. 1, p.
1--19, 2015.}

\refitem{MACAMBIRA, Indira Quaresma; ATHAYDE, Lígia S. Reprovação na
disciplina cálculo nos cursos de engenharia: análise de dados e métodos
minimizadores. \emph{In}: CONGRESSO BRASILEIRO DE EDUCAÇÃO EM ENGENHARIA
(COBENGE), 42., 2014. \textbf{Anais [\ldots]}. Brasília: ABENGE, 2014.}

\refitem{MÉXICO. Secretaría de Educación Pública (DGB/SEMS). \textbf{Marco
Curricular Común de la Educación Media Superior --- Pensamiento
Matemático}. Cidade do México: SEP, 2023.}

\refitem{MULLIS, Ina V. S. \emph{et al.} \textbf{TIMSS 2015 international
results in mathematics}. Chestnut Hill, MA: TIMSS \& PIRLS International
Study Center, Boston College, 2016.}

\refitem{OECD. \textbf{PISA 2022 results (Volume I)}: the state of
learning and equity in education. Paris: OECD Publishing, 2023.}

\refitem{PAQUISTÃO. Federal Board of Intermediate and Secondary Education
(FBISE). \textbf{Scheme of Studies --- Mathematics HSSC-II}. Islamabad:
FBISE, 2023.}

\refitem{PIRES, Célia Maria Carolino; GODOY, Elenilton Vieira; SILVA,
Márcio Antonio da; SANTOS, Vinício de Macedo. O currículo de Matemática em
revista: um editorial. \textbf{Bolema: Boletim de Educação Matemática},
Rio Claro, v. 28, n. 49, p. 485--493, 2014.}

\refitem{PNUD --- Programa das Nações Unidas para o Desenvolvimento.
\textbf{Relatório de Desenvolvimento Humano 2023/2024}. Nova York: PNUD,
2024.}

\refitem{REINO UNIDO. Department for Education. \textbf{Mathematics: AS
and A level content}. Londres: DfE, 2016. Ref. DFE-00706-2014.}

\refitem{REZENDE, Wanderley Moura. \textbf{O ensino de Cálculo}:
dificuldades de natureza epistemológica. 2003. Tese (Doutorado em
Educação) --- Faculdade de Educação, Universidade de São Paulo, São Paulo,
2003.}

\refitem{RÚSSIA. Federal Institute of Pedagogical Measurements (FIPI).
\textbf{Exame Estatal Unificado (EGE) --- Matemática, nível Profile}:
codificador de conteúdos. Moscou: FIPI, 2024.}

\refitem{SEAB --- Singapore Examinations and Assessment Board.
\textbf{Mathematics: Singapore-Cambridge GCE A-Level Higher 2 (Syllabus
9758)}. Singapura: SEAB; Ministry of Education, 2025.}

\refitem{SFARD, Anna. On the dual nature of mathematical conceptions:
reflections on processes and objects as different sides of the same coin.
\textbf{Educational Studies in Mathematics}, v. 22, n. 1, p. 1--36, 1991.}

\refitem{SILVA, Circe Mary Silva da. Positivismo, ensino secundário e
cálculo diferencial e integral. \textbf{Revista História da Educação},
v. 27, e128201, 2023. DOI: 10.1590/2236-3459/128201.}

\refitem{SOARES, Flávia. \textbf{Movimento da Matemática Moderna no
Brasil}: avanço ou retrocesso? 2001. Dissertação (Mestrado) --- Pontifícia
Universidade Católica do Rio de Janeiro, Rio de Janeiro, 2001.}

\refitem{SOCIEDADE BRASILEIRA DE EDUCAÇÃO MATEMÁTICA (SBEM). O Cálculo
Diferencial e Integral e o Ensino Médio. \emph{In}: ENCONTRO NACIONAL DE
EDUCAÇÃO MATEMÁTICA (ENEM), 9., 2007, Belo Horizonte. \textbf{Anais
[\ldots]}. Recife: SBEM, 2007.}

\refitem{TALL, David; VINNER, Shlomo. Concept image and concept definition in
mathematics with particular reference to limits and continuity.
\textbf{Educational Studies in Mathematics}, v. 12, n. 2, p. 151--169, 1981.}

\refitem{TURQUIA. Millî Eğitim Bakanlığı (MEB). \textbf{Ortaöğretim
Matematik Dersi Öğretim Programı}. Ancara: MEB, 2018.}

\refitem{VIETNÃ. Bo Giao duc va Dao tao (MOET). \textbf{Chuong trinh Giao
duc Pho thong mon Toan} (Circular 32/2018/TT-BGDDT). Hanói: MOET, 2018.}

\refitem{VYGOTSKY, Lev S. \textbf{Mind in society}: the development of
higher psychological processes. Cambridge, MA: Harvard University Press,
1978.}

\refitem{WAEC --- West African Examinations Council. \textbf{Further
Mathematics Syllabus (WASSCE/SSCE)}. Lagos: WAEC, 2023.}

\refitem{WOOD, David; BRUNER, Jerome S.; ROSS, Gail. The role of tutoring
in problem solving. \textbf{Journal of Child Psychology and Psychiatry},
v. 17, n. 2, p. 89--100, 1976.}

\endgroup

\end{document}
